\subsection{Scenarios}
\subsubsection*{User publishes a path}
\begin{enumerate}
    \item The user accesses the BBP system and opens the “My Routes” section.
    \item The system displays the list of paths associated with the user, indicating for each one, among other properties, its status (private or public).
    \item The user selects a private route and, using the appropriate button, requests its publication.
    \item Upon receiving the publication request, the system asks the user to confirm additional information related to the route:
        \begin{itemize}
            \item In the case of manually created routes, the user can review and, if necessary, modify the additional information previously entered during creation.
            \item In the case of automatically generated routes, the user must confirm or correct the acquired information (if not already done) and may add further details manually.
        \end{itemize}
    \item Once the review is complete, the user confirms the publication of the route.
    \item The system verifies the data consistency, assigns a unique identifier to the route, and makes it available to the community within the search and browsing functionality.
    \item The system notifies the user of the successful publication and updates the route status from private to public.
\end{enumerate}
\subsubsection*{User modifies the additional information related to a route}
\begin{enumerate}[label=\arabic*.]
    \item The user accesses the Best Bike Paths (BBP) platform and opens the “My Routes” section.
    \item The system displays the list of routes associated with the user, indicating for each one, among other properties, its status (private or public).
    \item The user selects a route and opens the section dedicated to additional information.
    \item The system displays the current information, including:
        \begin{itemize}
            \item The condition of the roads that make up the route;
            \item The presence of obstacles or damaged sections;
            \item The origin of each piece of data (manually entered or automatically acquired).
        \end{itemize}

    \item The user can modify, correct, or add new information:
        \begin{itemize}
            \item If the route was created manually, the user can freely update existing data.
            \item If the route was created automatically, the user can review the information acquired by the system and/or add new data manually; in any case, they must confirm or correct the automatically acquired information before publication.

        \end{itemize}
    \item The system saves the changes and updates the database.




\end{enumerate}
